% !TEX root = <name of your master file>
%
% Chapter: Strange Things in Classical Mechanics
%
% This file is written to be \input{} or \include{}d into a master
% document using the book (or report) class. It assumes the master
% preamble loads at least:
%   \usepackage{amsmath,amssymb,amsthm}
%   \usepackage{mathtools}
%   \usepackage[margin=1in]{geometry}   % optional
% No other special packages are required.

\chapter{Strange Things in Classical Mechanics}
\label{chap:strange-classical-mechanics}

\section{Introduction}
\label{sec:strange-cm-intro}

Classical mechanics is usually presented as the most settled part of
physics: three laws, a variational principle, and a semester of
worked examples. By the time a student reaches quantum theory or
general relativity, mechanics has already receded into the
background as a body of technique rather than a live subject. This
impression is not wrong, exactly, but it is incomplete. If one
presses on a few of the concepts that the standard course treats as
self-evident --- \emph{force}, \emph{constraint}, \emph{system} ---
the apparatus turns out to rest on assumptions that are easy to
state and easy to violate. What happens is instructive: relaxing a
single tacit assumption does not merely complicate the bookkeeping,
it can change the physics outright, produce equations that look
paradoxical on first encounter, or reveal that two "obviously
equivalent" formulations of mechanics are not equivalent at all.

This chapter collects four such cases.
\begin{itemize}
  \item Section~\ref{sec:hertz} looks at Heinrich Hertz's attempt to
  do mechanics \emph{without} taking force as a primitive concept at
  all, replacing it by pure geometry and concealed degrees of
  freedom.
  \item Section~\ref{sec:nonholonomic} looks at constraints whose
  defining equations cannot be integrated, and shows that for such
  constraints the two most natural-looking generalizations of
  Hamilton's principle give \emph{different, both self-consistent,
  but physically inequivalent} theories --- only one of which
  matches experiment.
  \item Section~\ref{sec:shape} looks at bodies that can change
  their own shape, and shows how a falling cat can reorient itself,
  and a bacterium can swim, without ever violating the conservation
  of angular momentum or momentum.
  \item Section~\ref{sec:open} looks at systems that are not closed
  --- that exchange mass or energy with an unmodelled environment
  --- and shows how "extra" force terms and outright pathologies
  (runaway solutions, apparent violations of causality) can appear
  the moment a system is treated as open.
\end{itemize}
A closing remark (Section~\ref{sec:common-thread}) draws the four
threads together.

\section{Hertz's Mechanics Without Force}
\label{sec:hertz}

\subsection{The problem Hertz set himself}

By the 1890s, several physicists --- Kirchhoff and Mach among them
--- had grown uneasy with the status of \emph{force} as a
fundamental concept of mechanics. Force is defined operationally
through $F=ma$, but that same equation is routinely used to
\emph{introduce} force, which makes the law read uncomfortably like
a definition dressed up as an empirical statement. In his
posthumously published \emph{Die Prinzipien der Mechanik in neuem
Zusammenhange dargestellt} (1894), Heinrich Hertz proposed to remove
the discomfort by removing force altogether: he set out to build the
whole of mechanics from just three primitive notions --- space,
time, and mass --- with no independent concept of force at
all~\cite{Hertz1894,Lutzen2005}.

\subsection{The principle of the straightest path}

Consider $N$ mass points with positions $\mathbf x_i\in\mathbb R^3$
and masses $m_i$, and equip the $3N$-dimensional configuration space
$Q$ with the natural mass-weighted (kinetic-energy) metric,
\begin{equation}
  ds^2 \;=\; \sum_{i=1}^N m_i\, d\mathbf x_i\cdot d\mathbf x_i,
  \qquad
  T \;=\; \frac12\left(\frac{ds}{dt}\right)^{\!2}.
  \label{eq:hertz-metric}
\end{equation}
Hertz's masses are never free in the naive sense: they are always
tied together by \emph{rigid, workless connections} (idealized
mechanical linkages that transmit force but do no work), which
restrict the accessible configurations to some submanifold
$\mathcal C\subset Q$. His fundamental postulate --- the
\emph{principle of the straightest path}
(\textit{Prinzip der geradesten Bahn}) --- is then simply:

\begin{quote}
\emph{A system free of "force" moves along the straightest possible
path compatible with its constraints, at constant speed.}
\end{quote}

"Straightest possible" is to be read literally, in the geometric
sense: the representative point traces a \emph{geodesic} of the
metric~\eqref{eq:hertz-metric} restricted to $\mathcal C$, i.e.\ a
curve whose acceleration (in the induced metric) is purely normal to
$\mathcal C$ and is entirely accounted for by the constraint
reaction. Formally, if the constraints are written $g_a(\mathbf
x)=0$, $a=1,\dots,k$, the equations of motion are
\begin{equation}
  m_i\,\ddot{\mathbf x}_i \;=\; \sum_a \lambda_a\,
  \frac{\partial g_a}{\partial \mathbf x_i},
  \label{eq:hertz-eom}
\end{equation}
with the multipliers $\lambda_a$ fixed by demanding $g_a(\mathbf
x(t))\equiv 0$. There is nothing on the right-hand side of
\eqref{eq:hertz-eom} except constraint reactions: no applied-force
term appears anywhere, by construction. Constant speed
($|ds/dt|=\mathrm{const}$, i.e.\ constant kinetic energy) follows
automatically, since ideal constraints do no work.

\subsection{Hidden masses, or where "force" goes}

Of course, ordinary bodies \emph{do} appear to exert forces on one
another --- gravity, contact forces, and so on all look like genuine
interactions between visible masses that are not rigidly connected
to anything. Hertz's resolution was to posit \emph{hidden masses}:
invisible mass points, rigidly linked to the visible ones (and, in
his more speculative moments, identified with a mechanical
\ae{}ther), whose only role is to curve the constraint submanifold
$\mathcal C$ in just the right way. If the full configuration space
splits as $Q = Q_{\mathrm{vis}}\times Q_{\mathrm{hid}}$ and the
\emph{complete} system --- visible plus hidden --- moves along a
geodesic of $\mathcal C$, then projecting that geodesic motion down
onto $Q_{\mathrm{vis}}$ alone generically produces a curve that is
\emph{not} geodesic in the visible sector by itself: extra terms
appear in its equation of motion that are, from the point of view of
an observer with access only to $Q_{\mathrm{vis}}$, indistinguishable
from the effect of an applied force~\cite{Hertz1894,Lutzen2005}.

This is a genuinely striking claim: on Hertz's account, force is not
absent from experience, only absent from the fundamental ontology.
Every observed force law would, in principle, be explained by some
specific (if permanently unobservable) arrangement of hidden mass
and rigid linkage. The same basic move --- explain an apparent force
as the shadow, on a submanifold, of pure geometric (geodesic) motion
in a larger space --- reappears twice in twentieth-century physics
in far more successful forms: in Kaluza--Klein theory, where the
electromagnetic force on a charged particle is recovered as geodesic
motion in a five-dimensional space, and, far more importantly, in
general relativity, where gravity itself --- which Newton had to
treat as a force acting at a distance --- becomes geodesic motion of
a free particle through curved four-dimensional spacetime, with no
"force" required at all. Hertz did not anticipate curved spacetime;
his geometry is Euclidean and the curvature is manufactured entirely
by hidden extra coordinates. But the structural resemblance to
Einstein's later, physically successful geometrization of gravity is
close enough that Hertz's book is still read today chiefly for that
reason, rather than as a working theory in its own
right~\cite{Lutzen2005}.

\subsection{A second, more conservative geometrization}

It is worth noting that ordinary, force-\emph{having} mechanics can
also be written as a geodesic principle, without any hidden masses
at all, via Jacobi's form of the principle of least action. For a
conservative system with potential $V(\mathbf x)$ and fixed total
energy $E$, the true trajectories are exactly the geodesics of the
\emph{Jacobi metric}
\begin{equation}
  ds_J^2 \;=\; 2\bigl(E-V(\mathbf x)\bigr)\, ds^2 ,
  \label{eq:jacobi-metric}
\end{equation}
a conformal rescaling of the mass metric~\eqref{eq:hertz-metric} by
the local kinetic energy. So there are, in the end, two quite
different ways to eliminate force from the description of a system
governed by a potential: enlarge the configuration space with hidden
coordinates and keep the metric flat (Hertz), or keep the visible
configuration space as it is and let the metric itself become curved
and energy-dependent (Jacobi). Both reproduce exactly the same
observable trajectories. Neither is more "true" than the other; they
are, in Hertz's own vocabulary, different \emph{images} of the same
phenomena. That a concept as apparently basic as force can be made
to evaporate --- twice over, in two structurally different ways ---
without changing a single observable prediction is the real content
of this section: it shows "force" to be a feature of a chosen
representation of mechanics, not a fact about the world that every
correct representation must contain.

Hertz's own programme did not survive contact with its critics.
Almost every contemporary reviewer, Boltzmann included, objected
that positing a fresh, permanently unobservable mechanism of hidden
masses for every distinct force law (gravitation, electromagnetism,
\ldots) bought elegance at the price of testability, and the
research programme was essentially abandoned within a generation. It
survives today less as working physics than as an early, and
unusually explicit, demonstration that mechanics has more than one
internally consistent axiomatic foundation --- a lesson later put to
far more productive use.

\section{Nonholonomic Constraints: When Differentials Refuse to
  Integrate}
\label{sec:nonholonomic}

\subsection{Holonomic versus nonholonomic}

A constraint of the form $f(q,t)=0$, where $q=(q^1,\dots,q^n)$ are
generalized coordinates, is called \emph{holonomic}: it can be
solved, at least locally, to eliminate one coordinate outright, and
it reduces the dimension of the configuration space itself. A more
general class of constraints is linear in the velocities,
\begin{equation}
  \sum_{i=1}^n a_i(q,t)\,\dot q^i + a_t(q,t) \;=\;0,
  \label{eq:pfaffian}
\end{equation}
a so-called Pfaffian constraint. If the left-hand side of
\eqref{eq:pfaffian} happens to be the total time derivative of some
function $f(q,t)$, the constraint is secretly holonomic in disguise.
If it is not --- if no such $f$ exists --- the constraint is called
\emph{nonholonomic}: it restricts the \emph{velocities} available at
each point of $Q$, without restricting which \emph{configurations}
are reachable.

The canonical example is a coin (or disk) of radius $R$ rolling
without slipping on a plane, with contact point $(x,y)$, heading
angle $\phi$, and roll angle $\theta$ as coordinates. The
rolling condition reads
\begin{equation}
  \dot x = R\,\dot\theta\cos\phi ,\qquad
  \dot y = R\,\dot\theta\sin\phi ,
  \label{eq:rolling-coin}
\end{equation}
which is not integrable: it constrains the instantaneous velocity to
a two-dimensional subspace of the four-dimensional tangent space at
each point, and yet, by alternating steering and rolling, the coin
can be driven from \emph{any} configuration $(x,y,\phi)$ to any
other. The instantaneous constraint is genuine, but it places no
restriction whatsoever on the reachable set of configurations --- an
instance of what control theorists call bracket-generating, or
Chow--Rashevski\u{i}, controllability. This already runs against a
piece of intuition carried over uncritically from the holonomic
case, where a nontrivial velocity constraint always shrinks the
accessible configuration space too.

A second standard example, the \emph{Chaplygin sleigh}, makes an
even sharper point. It is a rigid body sliding frictionlessly on a
plane, with a knife edge that permits motion only along its own
length (no sideways slip). The system conserves energy exactly ---
there is no friction anywhere in the model --- and yet its reduced
dynamics on the velocity variables is attracted to a fixed
equilibrium: left alone, the sleigh asymptotically settles into
straight-line motion with no rotation, knife edge trailing the
center of mass~\cite{Chaplygin1911,ChaplyginSleighWiki}. This
happens because the reduced equations of motion on the nonholonomic
constraint surface, unlike those of an ordinary holonomic mechanical
system, do not in general preserve phase-space (Liouville) volume,
even though they do preserve energy. A perfectly conservative
mechanical system can therefore display genuinely dissipative-looking
long-time behaviour --- attracting fixed points, basins of
attraction --- with no friction anywhere in the equations. Energy
conservation, by itself, is not enough to rule out the qualitative
signatures usually associated with dissipation once constraints stop
being holonomic.

\subsection{Two variational principles, two different theories}

For holonomic constraints, two routes to the equations of motion ---
D'Alembert's principle (constraint forces do no virtual work, for
virtual displacements compatible with the instantaneous constraint)
and Hamilton's principle (the true path extremizes the action among
\emph{all} paths satisfying the constraint) --- are provably
equivalent. For nonholonomic (non-integrable, velocity-linear)
constraints, this equivalence simply fails, and the failure is not a
minor technicality.

There are two natural-looking ways to generalize the variational
treatment to a system with a nonholonomic constraint of the
form~\eqref{eq:pfaffian}.

\begin{description}
  \item[Nonholonomic (Lagrange--d'Alembert, or Chetaev) mechanics.]
  Impose the constraint directly on virtual displacements at each
  instant, exactly as in the holonomic case, and only afterwards
  derive equations of motion:
  \begin{equation}
    \frac{d}{dt}\frac{\partial L}{\partial \dot q^i}
    - \frac{\partial L}{\partial q^i}
    \;=\; \sum_a \lambda_a\, a_i^a(q),
    \qquad
    \sum_i a_i^a(q)\,\dot q^i + a_t^a(q) = 0 .
    \label{eq:lagrange-dalembert}
  \end{equation}
  \item[Vakonomic mechanics.] (The name --- "variational axiomatic
  kind" --- was coined by Kozlov~\cite{Kozlov1985,VershikFaddeev}.)
  Instead, treat the nonholonomic constraint as a genuine constraint
  on the space of \emph{paths} from the outset, as one would an
  isoperimetric constraint in the calculus of variations, and only
  then extremize the action over that constrained set of paths.
\end{description}

For holonomic constraints these two prescriptions coincide. For
genuinely nonholonomic constraints they generically do not: it has
been shown, for instance, that for the rolling coin on an inclined
plane no nontrivial vakonomic solution reproduces the
Lagrange--d'Alembert solution at
all~\cite{Zampieri2005,ZenkovEtAl2021}. This is not merely a
mathematical curiosity to be settled by taste: Lewis and Murray
built and instrumented a physical system --- a ball rolling without
slipping on a rotating table --- specifically to adjudicate between
the two theories, and found that the observed motion matches the
Lagrange--d'Alembert (nonholonomic) prediction; they were unable to
produce vakonomic simulations resembling the experimental data at
all~\cite{LewisMurray1995}. The variational principle that looks like
the most direct generalization of Hamilton's principle to velocity
constraints is, for constraints realized by an actual rolling or
sliding contact, simply the wrong physical theory.

This does not make vakonomic mechanics useless --- it correctly
describes optimal-control and sub-Riemannian-geodesic problems, where
the "constraint" genuinely is a restriction on the class of paths
being optimized over, rather than a constraint enforced
instant-by-instant by a physical reaction force. The moral is
subtler and more interesting than "one theory is right and the other
wrong": the bare constraint equation~\eqref{eq:pfaffian}, by itself,
does not determine the dynamics. The correct equations of motion
depend on how the constraint is physically \emph{realized} ---
information that is invisible at the level of the constraint
equation alone, and that a first course in mechanics, built entirely
on holonomic examples where this ambiguity cannot arise, gives no
reason to expect.

\section{Bodies That Change Shape: Locomotion Without Anything to
  Push Against}
\label{sec:shape}

\subsection{The falling cat}

Drop a cat, upside down, from rest. In free fall there is no
external torque about the center of mass, so angular momentum is
conserved --- and it starts, and remains, at zero. Naive intuition
says that a body with identically zero angular momentum cannot
rotate. And yet the cat lands on its feet, having reoriented itself
by roughly $180^\circ$. The naive inference silently assumes a
\emph{rigid} body; it simply does not apply to a body that can
deform itself internally, and the cat's righting maneuver --- a
sequence of twists of its front and back halves relative to one
another --- is a real, controlled deformation, not a violation of
any conservation law.

\subsection{Shape space and the gauge structure of deformation}

The clean way to see why this works, due to Shapere and
Wilczek~\cite{ShapereWilczek1989} and to Montgomery's later "gauge
theory of the falling cat"~\cite{Montgomery1993}, is to separate a
deformable body's configuration into its \emph{shape} $s$ (the
internal, orientation-independent degrees of freedom --- relative
joint angles, say) and an overall orientation $R\in SO(3)$ that fixes
that shape in space. Locally, the full configuration space is a
principal fiber bundle with structure group $SO(3)$ over shape space,
and the requirement that total angular momentum stay identically
zero throughout the motion is a linear relation between the shape
velocity $\dot s$ and the angular velocity $\omega$ of the body,
\begin{equation}
  \omega \;=\; -\,\mathcal A(s)\,\dot s ,
  \label{eq:connection}
\end{equation}
where $\mathcal A(s)$ is built from the shape-dependent inertia
tensor and its couplings. This is precisely the data of a connection
(a gauge potential) on the bundle: given any path $s(t)$ in shape
space, \eqref{eq:connection} lifts it uniquely to a path in the full
configuration space, exactly as a connection lifts a path in the
base of a fiber bundle to a horizontal path upstairs.

The consequence that matters here is that if the shape returns
exactly to where it started after a closed loop in shape space,
$s(T) = s(0)$, the orientation need \emph{not} return to where it
started. The net rotation accumulated around the loop is the
\emph{holonomy} of the connection~\eqref{eq:connection}, and for a
small loop enclosing area $\Sigma$ it is governed, Stokes'-theorem
style, by the curvature $\mathcal F = d\mathcal A + \mathcal
A\wedge\mathcal A$ of that connection:
\begin{equation}
  \Delta R \;\approx\; \exp\!\left[\oint_{\partial\Sigma}\mathcal
  A\right] \;\approx\; \exp\!\left[\int_{\Sigma}\mathcal F\right].
  \label{eq:holonomy}
\end{equation}
This is exactly the structure of a non-Abelian geometric (Berry,
Wilczek--Zee) phase from quantum mechanics, and of Aharonov--Bohm
holonomy in gauge theory --- reproduced in full by an entirely
classical, entirely mechanical system. It is also, viewed from the
vantage point of Section~\ref{sec:nonholonomic}, nothing but another
instance of a nonholonomic constraint: "total angular momentum
$\equiv 0$" is a Pfaffian relation between $\dot s$ and $\omega$ that
fails to be integrable to any function of shape and orientation
alone \emph{exactly when} the connection~\eqref{eq:connection} has
nonzero curvature. The falling-cat effect and the rolling coin are,
formally, the same phenomenon in different clothing.

\subsection{Swimming without pushing}

The same structure governs self-propulsion in a highly viscous
fluid, where inertia is negligible and net displacement over one
cycle of body deformation is possible only if the cycle of shapes is
genuinely non-reciprocal. This is Purcell's scallop theorem: a
scallop that opens and closes through the \emph{same} sequence of
shapes, forwards then backwards, achieves zero net displacement,
because time-reversing the (inertia-free) equations of motion
exactly reverses the trajectory. Net swimming is again a
holonomy/geometric-phase effect in an appropriately defined shape
bundle~\cite{ShapereWilczekSwim1989}.

Perhaps the most striking illustration of the underlying idea is Jack
Wisdom's result that a quasi-rigid body executing a cyclic sequence
of internal shape changes in a curved \emph{spacetime} --- as
prescribed by general relativity --- can achieve a net translation
even in complete vacuum, with nothing to push against at
all~\cite{Wisdom2003}. The effect scales with the local spacetime
curvature and with the extent of the body's internal motion; it is
minuscule for any realistic body and is not a practical propulsion
mechanism, but as a matter of principle it is "swimming" through
spacetime using exactly the same holonomy mechanism as the falling
cat, now realized by the genuine nonlinearity of general relativity
rather than by a gauge potential on a shape bundle constructed by
hand. Once again, an outcome that sounds like it must violate
momentum conservation turns out to be perfectly consistent with
conservation laws, provided one tracks correctly what "conservation"
does and does not forbid for a body that is allowed to change shape.

\section{Open Systems: Mechanics Without a Fixed Body of Matter}
\label{sec:open}

\subsection{Variable mass and the Meshchersky equation}

Ordinary Newtonian mechanics implicitly assumes a fixed set of
particles. Rockets, raindrops growing by accretion, chains lifted
off a pile, and avalanches all violate that assumption: the mass $m$
of the body under study is itself a function of time. The correct
generalization of Newton's second law to such an \emph{open system}
--- worked out around the turn of the twentieth century by
I.\,V.\ Meshchersky --- is
\begin{equation}
  m\,\frac{d\mathbf v}{dt} \;=\; \mathbf F_{\mathrm{ext}}
  + \mathbf u\,\frac{dm}{dt} ,
  \label{eq:meshchersky}
\end{equation}
where $\mathbf u$ is the velocity of the incoming or outgoing mass
\emph{relative to the body itself}, not relative to the lab frame.
Setting $\mathbf F_{\mathrm{ext}}=0$ and integrating
\eqref{eq:meshchersky} with $\mathbf u$ constant and antiparallel to
the direction of travel yields the Tsiolkovsky rocket equation,
\begin{equation}
  v(t) - v(0) \;=\; u\,\ln\frac{m(0)}{m(t)} .
  \label{eq:tsiolkovsky}
\end{equation}
The resemblance to Section~\ref{sec:hertz} is not superficial. The
term $\mathbf u\,\dot m$ in \eqref{eq:meshchersky} looks, from the
point of view of the rocket body alone, exactly like an applied
force --- "thrust" --- and is routinely called one. But viewed as
part of the larger \emph{closed} system consisting of the rocket
\emph{together with} every parcel of exhaust it has ever ejected,
whose total momentum is exactly conserved, there is no force there
at all: only the ordinary transfer of momentum between two pieces of
one conserved whole. What looks like a force is, once again, an
artifact of insisting on an equation of motion for part of a system
rather than for the whole of it.

This is also a genuine pedagogical trap, worth flagging explicitly.
It is tempting to write $\mathbf F = d(m\mathbf v)/dt = m\dot{\mathbf
v} + \mathbf v\,\dot m$ and treat this directly as the equation of
motion of the open body. That expression is correct as a statement
about the rate of change of the momentum of the body \emph{alone},
but it is not, in general, equivalent to \eqref{eq:meshchersky} and
gives the wrong dynamics unless the accreted or ejected mass happens
to leave (or join) with zero velocity relative to the body, i.e.\
$\mathbf u \equiv \mathbf v$ --- a special case, not the generic one,
and not the one realized by a rocket engine, whose whole purpose is
to eject mass at high relative velocity.

\subsection{Dissipative systems: putting friction back into a
  variational principle}

A second, quite different kind of "openness" arises whenever a
system exchanges energy with an environment that is not explicitly
modelled --- friction, drag, radiation reaction. None of these
forces derive from a potential in the ordinary sense, so the
Lagrangian and Hamiltonian formalisms, built for closed and
energy-conserving systems, do not apply directly. Two classical
devices repair this.

The first is Rayleigh's dissipation function $\mathcal
F_{\mathrm{diss}}(\dot q)$, quadratic in the generalized velocities
for linear damping, which modifies the Euler--Lagrange equations to
\begin{equation}
  \frac{d}{dt}\frac{\partial L}{\partial \dot q^i}
  - \frac{\partial L}{\partial q^i}
  \;=\; -\,\frac{\partial \mathcal F_{\mathrm{diss}}}{\partial \dot
  q^i}.
  \label{eq:rayleigh}
\end{equation}
This correctly reproduces the (irreversible) equations of motion,
but it is a hybrid: $\mathcal F_{\mathrm{diss}}$ is not itself part
of a genuine, self-contained stationary-action principle for $L$
alone.

The second, due to Bateman (1931), is more radical and more
illuminating: a fully conservative, symplectic Lagrangian \emph{can}
be written for a damped oscillator, but only at the price of
doubling the number of degrees of freedom. Introduce an auxiliary
"mirror" coordinate $y$ alongside the physical coordinate $x$,
\begin{equation}
  L_{\mathrm{dual}}(x,\dot x,y,\dot y) \;=\;
  m\,\dot x\,\dot y \;+\; \frac{\gamma}{2}\bigl(x\dot y - \dot x
  y\bigr) \;-\; k\,xy ,
  \label{eq:bateman}
\end{equation}
whose Euler--Lagrange equations are precisely the damped oscillator,
$\ddot x + (\gamma/m)\dot x + \omega_0^2 x = 0$, for the physical
coordinate $x$, paired with a time-reversed, exponentially
\emph{amplified} equation for $y$~\cite{Bateman1931}. The energy the
physical oscillator loses is exactly the energy its mirror partner
gains, so the doubled system is conservative even though either half
alone is not. This is a small, entirely classical and elementary
demonstration of an idea that later became central to the theory of
open quantum systems: dissipation and irreversibility in a subsystem
are not exceptions to an underlying conservative dynamics, but the
generic appearance of a conservative dynamics once part of the full
system --- here, quite literally, the mirror half --- has been
discarded from the description.

\subsection{When "open" turns pathological: radiation reaction}

Not every attempt to write a closed-looking equation for an open
system ends benignly. An accelerating point charge radiates, and by
momentum conservation it must experience a recoil, or self-force, in
consequence. Treating the charge \emph{in isolation} --- an open
system in the strict sense, since the charge together with its own
radiated field is really a system with infinitely many degrees of
freedom, formally "integrated out" down to a single effective force
term --- leads to the (nonrelativistic) Abraham--Lorentz equation,
\begin{equation}
  m\,\ddot{\mathbf x} \;=\; \mathbf F_{\mathrm{ext}} \;+\;
  m\tau\,\dddot{\mathbf x} ,
  \qquad
  \tau \;=\; \frac{q^2}{6\pi\varepsilon_0\,m\,c^3}
  \;=\; \frac{2}{3}\,\frac{r_e}{c},
  \label{eq:abraham-lorentz}
\end{equation}
where $r_e$ is the classical electron radius. Because
\eqref{eq:abraham-lorentz} is third order in position, it requires
more initial data than position and velocity to determine a unique
solution, and generic initial data produce \emph{runaway} solutions
in which the self-acceleration grows without bound even with
$\mathbf F_{\mathrm{ext}}=0$. Suppressing the runaway by hand forces
the opposite pathology: the charge must begin accelerating
\emph{before} an external force is switched on --- a small, short-lived,
but genuinely acausal-looking pre-acceleration. (In practice, both
pathologies signal that \eqref{eq:abraham-lorentz} is being pushed
outside its domain of validity, on timescales comparable to $\tau
\sim 10^{-24}\,\mathrm s$; treated properly, as an effective equation
valid only for slowly varying external forces --- via, e.g., the
Landau--Lifshitz reduction of order --- the pathologies disappear.)
The lesson is a cautionary complement to the rest of this section:
"integrating out" the unmodelled part of an open system is not
always harmless bookkeeping. Sometimes it is, as with the Meshchersky
equation; sometimes, as here, it produces an equation that is simply
malformed unless handled with real care.

\section{A Common Thread}
\label{sec:common-thread}

Each of the four cases in this chapter chips away at one piece of
the implicit picture underlying an introductory treatment of
mechanics: that force is a well-defined primitive; that Hamilton's
principle generalizes to velocity constraints in the one obvious
way; that a body with zero momentum or angular momentum cannot move
or turn; and that a mechanical system is a fixed, closed collection
of particles amenable to an ordinary Lagrangian. None of these
assumptions is stated explicitly in a first course, precisely
because within the scope of the standard worked examples they never
fail. They are, in that sense, exactly the kind of tacit
identification of a convenient special case with the general rule
that produces apparent paradoxes once the special case is left
behind. What is notable is how much genuine physics --- not merely
mathematical bookkeeping --- follows from tracking down each
assumption and relaxing it: gauge structures governing locomotion,
an experimentally decided answer to a real question about which
variational principle is physically correct, working spacecraft, and
the classical seed of the modern theory of open quantum systems. A
subject can be taught as settled and still, on inspection, be full of
open conceptual work.

\begin{thebibliography}{99}

\bibitem{Hertz1894}
H.~Hertz, \emph{Die Prinzipien der Mechanik in neuem Zusammenhange
  dargestellt} (Barth, Leipzig, 1894); English translation
  \emph{The Principles of Mechanics Presented in a New Form},
  transl.\ D.~E.\ Jones and J.~T.\ Walley (Macmillan, London, 1899).

\bibitem{Lutzen2005}
J.~L\"utzen, \emph{Mechanistic Images in Geometric Form: Heinrich
  Hertz's `Principles of Mechanics'} (Oxford University Press,
  2005).

\bibitem{Chaplygin1911}
S.~A.~Chaplygin, ``On the theory of motion of nonholonomic
  systems. The reducing-multiplier theorem,'' \emph{Mat.\ Sbornik}
  \textbf{28}, 303--314 (1911) [in Russian].

\bibitem{ChaplyginSleighWiki}
See, e.g., A.~V.~Borisov and I.~S.~Mamaev, ``The dynamics of a
  Chaplygin sleigh,'' \emph{J.\ Appl.\ Math.\ Mech.} \textbf{73},
  156--161 (2009), for a modern treatment of the sleigh's
  volume-contracting, dissipative-looking reduced dynamics.

\bibitem{Kozlov1985}
V.~V.~Kozlov, ``On the integration theory of the equations of
  nonholonomic mechanics,'' \emph{Adv.\ Mech.} \textbf{8}, no.~3,
  85--107 (1985) [in Russian]; English translation,
  \emph{Regular and Chaotic Dynamics} \textbf{7}, 161--176 (2002).

\bibitem{VershikFaddeev}
V.~I.~Arnold, V.~V.~Kozlov, and A.~I.~Neishtadt,
  \emph{Mathematical Aspects of Classical and Celestial Mechanics},
  3rd ed.\ (Springer, Berlin, 2006), Ch.~1, \S4.

\bibitem{Zampieri2005}
G.~Zampieri, ``Nonholonomic versus vakonomic dynamics,''
  \emph{J.\ Differential Equations} \textbf{163}, 335--347 (2000).

\bibitem{ZenkovEtAl2021}
See, e.g., the discussion and further references in
  ``Complete inequivalence of nonholonomic and vakonomic mechanics:
  rolling coin on an inclined plane,'' \emph{Acta Mech.} (2021).

\bibitem{LewisMurray1995}
A.~D.~Lewis and R.~M.~Murray, ``Variational principles for
  constrained systems: theory and experiment,''
  \emph{Int.\ J.\ Non-Linear Mech.} \textbf{30}, 793--815 (1995).

\bibitem{ShapereWilczek1989}
A.~Shapere and F.~Wilczek, ``Gauge kinematics of deformable
  bodies,'' \emph{Am.\ J.\ Phys.} \textbf{57}, 514--518 (1989).

\bibitem{Montgomery1993}
R.~Montgomery, ``Gauge theory of the falling cat,''
  \emph{Fields Inst.\ Commun.} \textbf{1}, 193--218 (1993).

\bibitem{ShapereWilczekSwim1989}
A.~Shapere and F.~Wilczek, ``Geometry of self-propulsion at low
  Reynolds number,'' \emph{J.\ Fluid Mech.} \textbf{198}, 557--585
  (1989).

\bibitem{Wisdom2003}
J.~Wisdom, ``Swimming in spacetime: motion by cyclic changes in
  body shape,'' \emph{Science} \textbf{299}, 1865--1869 (2003).

\bibitem{Bateman1931}
H.~Bateman, ``On dissipative systems and related variational
  principles,'' \emph{Phys.\ Rev.} \textbf{38}, 815--819 (1931).

\end{thebibliography}
